% Snapshot of main.tex taken on 2026-11-25: same mechanism, false-name proof via the quadratic payment certificate (14 pages).
% Not part of the build. To compile it, copy it over ../main.tex in the repository root.
\documentclass[format=acmsmall, review=false]{acmart}
\usepackage{acm-ec-26}
\usepackage{booktabs} % For formal tables
\usepackage[ruled]{algorithm2e} % For algorithms
\usepackage{amsthm} % For theorem and remark environments
% amssymb is not loaded: acmart's newtxmath already provides its symbols,
% and loading both fails with "Command \Bbbk already defined".
\usepackage{bbm}
\newcommand{\E}{\mathbb{E}}
\newcommand{\one}{\mathbb{1}}
\renewcommand{\algorithmcfname}{ALGORITHM}
\SetAlFnt{\small}
\SetAlCapFnt{\small}
\SetAlCapNameFnt{\small}
\SetAlCapHSkip{0pt}
\IncMargin{-\parindent}
\newcommand{\OPT}{\mathrm{OPT}}
\newcommand{\SW}{\mathrm{SW}}
\newcommand{\eps}{\varepsilon}

\usepackage{xcolor}

\newcommand{\bruno}[1]{{\color{red}Bruno: #1}}
% Choose a citation style by commenting/uncommenting the appropriate line:
%\setcitestyle{acmnumeric}
\setcitestyle{authoryear}

% Title. Note the optional short title for running heads. In the interest of anonymization, please do not include any acknowledgements.
\title[On sybil-proof combinatorial auctions]{Roses are red, violets are blue, and combinatorial false-name-proof auctions can be efficient too}

% Anonymized submission.
\author{Bruno Mazorra}

% Abstract. Note that this must come before \maketitle.
\begin{abstract}
Iwasaki et al. \cite{iwasaki2010worst} show that, under anonymity, symmetry across goods, and independence of irrelevant goods (IIG), the optimal worst-case welfare approximation ratio of deterministic false-name-proof combinatorial auctions with single-minded bidders is $(m+1)/2$, where $m$ is the number of goods. In this paper, we make two contributions. First, we show that this bound continues to hold without the IIG assumption. Second, we show that allowing randomization substantially improves worst-case welfare. Specifically, we construct a randomized mechanism that is universally truthful, false-name-proof in expectation, runs in polynomial time, and achieves a worst-case welfare approximation ratio of $O(\sqrt{m}\log m)$. Our mechanism incurs only an additional $O(\log m)$ factor in welfare approximation compared to the classical greedy algorithm, showing that randomization makes false-name-proofness compatible with computational efficiency.
\end{abstract}

\begin{document}

% Title page for title and abstract only.
\begin{titlepage}
\maketitle
%\tableofcontents

\end{titlepage}
\section{Introduction}
\section{Related work}
\section{Model}
\label{sec:model}

We consider combinatorial auctions with a single seller, a set
$M=[m]$ of indivisible goods, and a finite set $N$ of bidders,
where $m\ge1$ and $n=|N|$. Bidder $i$ has a private type
$\theta_i=(S_i,v_i)$, consisting of a nonempty target
$S_i\subseteq M$ and a value $v_i\ge0$. With a slight abuse of
notation, we also use $v_i$ for its single-minded valuation
function: $v_i(R)=v_i\mathbf{1}\{S_i\subseteq R\}$ for every
bundle $R\subseteq M$. We denote the type space by $\Theta_M$.
Bidders are risk-neutral and have quasi-linear utilities:
receiving $R$ and paying $P$ gives utility $v_i(R)-P$.


A direct mechanism $\mathcal M=(X,P)$ maps a report profile
$\hat\theta=(\hat\theta_i)_{i\in N}$, with
$\hat\theta_i=(\hat S_i,b_i)\in\Theta_M$, to allocated bundles
$X_i(\hat\theta)\subseteq M$ and payments $P_i(\hat\theta)$
to the seller. The mechanism is defined for every finite set of
participating identities, including the empty set. Allocations
and payments may be randomized; all reports and identities are
fixed before the randomness is realized. Throughout, we require
\emph{ex-post feasibility}: for every report profile,
$X_i(\hat\theta)\cap X_j(\hat\theta)=\varnothing$ almost surely
for distinct $i,j$. Goods may remain unallocated. When no
confusion can arise, we drop the hat on reported targets and
write a report as $(S_i,b_i)$.

Write $u_i(\theta_i;\hat\theta)
=v_i(X_i(\hat\theta))-P_i(\hat\theta)$ for realized utility, and
\[
x_i(\hat\theta)
=\Pr[\hat S_i\subseteq X_i(\hat\theta)],
\qquad
\pi_i(\hat\theta)=\mathbb{E}[P_i(\hat\theta)]
\]
for service probabilities and expected payments. At a truthful
profile $\theta$, bidder $i$'s expected utility is
$U_i(\theta)=\mathbb{E}[u_i(\theta_i;\theta)]
=v_i x_i(\theta)-\pi_i(\theta)$.

We require \emph{individual rationality} (IR) in expectation,
meaning $U_i(\theta)\ge0$ for every bidder and type profile, and
\emph{no positive transfers} (NPT), meaning
$P_i(\hat\theta)\ge0$ almost surely for every report profile.
The expectation in IR is only over the mechanism's randomness,
with all other reports fixed, not over other bidders' types.
Thus $0\le\pi_i(\hat\theta)\le b_i x_i(\hat\theta)$.
For deterministic mechanisms, this is ex-post IR. For randomized
mechanisms, \emph{ex-post IR} additionally requires truthful
realized utility to be nonnegative almost surely.

\subsection{Truthfulness and false-name-proofness}

Let $E$ denote an arbitrary finite profile of external reports,
and write $(\hat\theta_i,E)$ for the profile obtained by adding
identity $i$ with report $\hat\theta_i$. A mechanism is
\emph{truthful} if, for every true type $\theta_i$, alternative
report $\hat\theta_i$, and external profile $E$,
\[
U_i(\theta_i,E)
\ge
\mathbb{E}\bigl[u_i(\theta_i;(\hat\theta_i,E))\bigr].
\]
For deterministic mechanisms, this is dominant-strategy
incentive compatibility (DSIC); for randomized mechanisms, it is
truthfulness in expectation. A deviation may change both the
target bundle and its value. A randomized mechanism is
\emph{universally truthful} if truthful reporting maximizes
realized utility for every realization of its randomness, that
is, if it is a distribution over DSIC mechanisms
\citep{dobzinski2006truthful}.

\emph{False-name-proofness} (FNP) \citep{yokoo2004effect} allows a
physical bidder to
replace its truthful report with any finite family of reports
$\hat\theta_A=(\hat\theta_j)_{j\in A}$ under distinct identities
$A$, its \emph{aliases}, not appearing in
$E$. The bidder values the union of their allocations and pays
the sum of their payments. We require
\[
U_i(\theta_i,E)
\ge
\mathbb{E}\!\left[
  v_i\!\left(\bigcup_{j\in A}X_j(\hat\theta_A,E)\right)
\right]
-\sum_{j\in A}\pi_j(\hat\theta_A,E)
\]
for every true type, external profile, and finite report family.
For randomized mechanisms, this is \emph{FNP in expectation};
for deterministic mechanisms, the expectation disappears.
The family may include the original identity, and different
identities may submit identical reports. Creating identities
is costless, and there is no bound on their number. We allow
$A=\varnothing$, representing nonparticipation, with empty unions
and sums interpreted as $\varnothing$ and zero. All submitted
identities remain in the profile and all assessed payments must
be paid: we consider ordinary FNP, without withdrawal. Taking
$|A|=1$ recovers truthfulness.

\subsection{Anonymity, symmetry, and irrelevant goods}

A mechanism is \emph{anonymous in utilities} if relabelling
bidders leaves their truthful expected utilities unchanged.
For a bijection $\tau:N\to N'$, define
$(\tau\theta)_{\tau(i)}=\theta_i$. Anonymity requires
$U_{\tau(i)}(\tau\theta)=U_i(\theta)$ for every profile, bidder,
and bijection. This definition applies to deterministic and
randomized mechanisms alike, and permits identity-based
tie-breaking when it leaves utilities unchanged.

Separately, a mechanism is \emph{symmetric across goods} if
$U_i(\sigma\theta)=U_i(\theta)$ for every bidder, profile, and
permutation $\sigma$ of $M$. Here $\sigma\theta$ replaces each
type $(S_j,v_j)$ by $(\sigma(S_j),v_j)$, so utilities on the
left are evaluated using the relabelled targets.

\begin{lemma}[Utility symmetry and critical prices]
\label{lem:utility-symmetry}
For a deterministic truthful mechanism satisfying IR and NPT,
anonymity in utilities is equivalent to anonymity of critical
prices. Likewise, symmetry across goods in utilities is
equivalent to symmetry across goods of critical prices.
\end{lemma}

\begin{proof}
Fix an identity $i$, an external profile $E$, and a target $S$.
DSIC makes service monotone in the reported value and gives
all served reports the same payment. IR and NPT imply that
unserved reports pay zero and that the winning payment equals
the critical price $c_{i,E}(S)$, the infimum of winning bids.
Consequently,
\[
U_i((S,v),E)=\max\{0,v-c_{i,E}(S)\}.
\]
If the critical price is infinite, this utility is identically
zero. At a finite critical price, winning and losing both give
zero utility, so tie-breaking does not affect the formula.
The critical price can therefore be recovered from utilities:
\[
c_{i,E}(S)
=\inf\{v\ge0:U_i((S,v),E)>0\},
\]
with $\inf\varnothing=+\infty$. Thus invariance of utilities
under either relabelling implies the corresponding invariance
of critical prices; the converse follows from the utility
formula. Under bidder anonymity, we may consequently write a
common menu $c_E(S)$ depending only on $S$ and the multiset
of external reports. Symmetry across goods further gives
$c_{\sigma E}(\sigma(S))=c_E(S)$.
\end{proof}

For comparison with the deterministic benchmark of
\cite[Definition~2]{iwasaki2010worst}, \emph{independence of
irrelevant goods} (IIG) requires that, for sufficiently large
values, a bidder winning the grand bundle continues to win the
enlarged grand bundle when a new good is added only to its
target, with all values and other targets unchanged.

\subsection{Worst-case welfare}

At a truthful profile $\theta$, realized welfare is
$\mathrm{SW}_{\mathcal M}(\theta)=\sum_{i\in N}v_i(X_i(\theta))$,
with expectation $\sum_{i\in N}v_i x_i(\theta)$, and optimal
welfare is
$\mathrm{OPT}(\theta)=\max_{A\text{ feasible}}\sum_{i\in N}v_i(A_i)$.
For a fixed number $m$ of goods, we define the
\emph{worst-case welfare fraction}\footnote{Its reciprocal is
the usual worst-case welfare approximation ratio, with
$1/0=+\infty$.} by the following infimum over all finite
populations and type profiles on $m$ goods, taking the
expectation over the mechanism's randomness at each fixed
profile:
\[
\rho(\mathcal M)
=
\inf_{\theta:\,\mathrm{OPT}(\theta)>0}
\frac{\mathbb{E}[\mathrm{SW}_{\mathcal M}(\theta)]}
     {\mathrm{OPT}(\theta)}.
\]
\section{Worst-case welfare of deterministic mechanisms}
\label{sec:deterministic}

Throughout this section, the mechanism is deterministic, feasible,
ex-post IR and NPT, DSIC, and ordinary FNP.  It is defined for
every finite population and admits a common bidder-anonymous critical-price
menu, equivariant under permutations of the goods
(Lemma~\ref{lem:utility-symmetry}); as there, we write $c_E(S)$ for this menu at
external profile $E$.  Thus bids strictly above $c_E(S)$ are served, bids
strictly below it are not served, and a winner pays $c_E(S)$; behavior
at equality is fixed by the tie rule.  We do not assume independence of
irrelevant goods.

If $\rho(\mathcal M)>0$, every critical price is finite.  Indeed, if
$c_E(S)=\infty$, then a bidder with target $S$ loses at every finite
value $v$.  The mechanism's welfare is at most $\OPT(E)$, whereas the new
optimum is at least $v$; sending $v\to\infty$ contradicts a positive ratio.
We use this observation without further comment.

The following inequality is the basic deterministic consequence of
false-name-proofness.

\begin{lemma}[False-name cover inequality]
\label{lem:det-cover}
Fix an external profile $E$ and a true target $S$.  Let $\hat\theta_A$ be a
finite false-name family with reports $(S_j,b_j)$, $j\in A$, and let
$D\subseteq A$ be the aliases that win.  If their allocations cover $S$, then
\begin{equation}
  c_E(S)
  \le
  \sum_{j\in D}
  c_{E\cup\hat\theta_{A\setminus\{j\}}}(S_j).
  \label{eq:det-cover}
\end{equation}
Every losing alias remains in the price context on the right-hand side.
\end{lemma}

\begin{proof}
Give the physical bidder a true value $v>c_E(S)$.  Truthful reporting
then gives utility $v-c_E(S)$.  Under the false-name profile, the winning
aliases cover $S$, so the physical value is again $v$.  Losing aliases
pay zero by IR and NPT, while each winner pays the critical price shown
in \eqref{eq:det-cover}.  False-name-proofness lets the common value $v$
cancel and gives the inequality.
\end{proof}

For $h>0$, write $H_h=(M,h)$ for the grand-bundle report of value $h$.
Goods equivariance gives a common clean singleton price
\[
  t(h)=c_{\{H_h\}}(\{g\}),
  \qquad g\in M.
\]

\begin{lemma}[Clean prices have a near-flat pair]
\label{lem:no-iig-flat}
For every $h>0$, $0<t(h)<\infty$.  Moreover, for every $\eps>0$ there exist
$h<h'$ in $[1,2]$ such that
\begin{equation}
  t(h')\le(1+\eps)t(h).
  \label{eq:near-flat}
\end{equation}
\end{lemma}

\begin{proof}
Finiteness was established above.  If $t(h)=0$, then every singleton bid
$\delta>0$ strictly wins against $H_h$.  Feasibility forces the grand-bundle
bidder to lose, so the welfare ratio on this two-bidder profile is at most
$\delta/h$.  Letting $\delta\downarrow0$ contradicts positivity.

For the second claim, fix an integer $\ell\ge1$ and consider the grid
points $1+s/\ell$, $s=0,\ldots,\ell$.  If every adjacent pair satisfied
$t(1+(s+1)/\ell)>(1+\eps)\,t(1+s/\ell)$, then
\[
  t(2)>(1+\eps)^\ell\, t(1).
\]
Because both endpoint prices are positive and finite, this fails for all
sufficiently large $\ell$.  One adjacent pair therefore satisfies
\eqref{eq:near-flat}.
\end{proof}

The next lemma is the replacement for IIG.  A collection of singleton bids
below half the clean price can neither lower another singleton's price nor
co-win with a higher singleton.

\begin{lemma}[Half-price persistence]
\label{lem:no-iig-half}
Fix $h>0$ and a good $g$, write $t=t(h)$, and choose $\lambda<t/2$.  Let $L$
be any family of distinct singleton reports, all of value $\lambda$, on goods
other than $g$.  Then
\begin{equation}
  c_{\{H_h\}\cup L}(\{g\})\ge t.
  \label{eq:no-iig-half}
\end{equation}
Moreover, if $(\{g\},y)$ wins against $\{H_h\}\cup L$ for some $y>\lambda$, no
report in $L$ co-wins.
\end{lemma}

\begin{proof}
We induct on $|L|$, proving the price bound and the no-co-winner statement
together.  Both are immediate when $L=\varnothing$.

Assume the result for every smaller family.  Suppose that $(\{g\},y)$ and a
low report $(\{g'\},\lambda)\in L$ both win.  Put
\[
  L'=L\setminus\{(\{g'\},\lambda)\},
  \qquad
  \gamma=c_{\{H_h,(\{g\},y)\}\cup L'}(\{g'\}).
\]
Because the low report wins, $\gamma\le\lambda<y$.  Raise its bid from
$\lambda$ to $y$.
Its own critical price remains $\gamma$, so it still wins.  Swapping goods
$g$ and $g'$ fixes $H_h$ and every report in $L'$; goods equivariance
therefore gives
\[
  c_{\{H_h,(\{g'\},y)\}\cup L'}(\{g\})=\gamma.
\]
Thus, against the smaller external profile $\{H_h\}\cup L'$, the two reports
$(\{g\},y)$ and $(\{g'\},y)$ both strictly win and pay $\gamma$ each.  Treat
them as two aliases of one bidder with true target $\{g\}$.  The cover
inequality and the induction hypothesis imply
\[
  t
  \le c_{\{H_h\}\cup L'}(\{g\})
  \le 2\gamma
  \le 2\lambda
  <t,
\]
a contradiction.  Hence a high singleton has no low co-winner.

Now choose
\[
  y>\max\{\lambda,c_{\{H_h\}\cup L}(\{g\})\}.
\]
The high singleton wins, and the statement just proved makes every report in
$L$ lose.  Regard the principal and the reports in $L$ as one false-name
family against $H_h$.  Only the principal pays, so
Lemma~\ref{lem:det-cover} gives
\[
  t=c_{\{H_h\}}(\{g\})
  \le c_{\{H_h\}\cup L}(\{g\}).
\]
This proves the price bound and completes the induction.
\end{proof}

\begin{theorem}[Deterministic upper bound without IIG]
\label{thm:no-iig-upper}
Every mechanism in the deterministic class above satisfies
\[
  \rho(\mathcal M)\le\frac{2}{m+1}.
\]
The bound already holds on single-minded star profiles.
\end{theorem}

\begin{proof}
The claim is immediate if $\rho(\mathcal M)=0$, so assume the ratio is
positive and $m\ge2$.  Fix $\eps>0$ and choose the near-flat pair $h<h'$ from
Lemma~\ref{lem:no-iig-flat}.  Write
\[
  t=t(h),
  \qquad
  t'=t(h'),
  \qquad
  \bar t=\max\{t,t'\}.
\]
Then $\bar t\le(1+\eps)t$.  Choose $a>\bar t$ and define the clean reverse
price
\[
  \eta(a)=c_{\{(\{1\},a)\}}(M).
\]
Because $a>t'$, the singleton $(\{1\},a)$ strictly wins against $H_{h'}$.
The grand-bundle report of value $h'$ therefore loses, so
\begin{equation}
  \eta(a)\ge h'>h.
  \label{eq:reverse-slack}
\end{equation}

Singleton auxiliaries cannot lower this grand-bundle price.  Indeed, for any
family $L$ of singleton reports, choose a grand-bundle bid above the finite
threshold $c_{\{(\{1\},a)\}\cup L}(M)$.  The grand report wins and every
singleton in $L$ loses.  Applying Lemma~\ref{lem:det-cover} to this principal
and the losing auxiliaries, against external report $(\{1\},a)$, gives
\begin{equation}
  c_{\{(\{1\},a)\}\cup L}(M)\ge\eta(a).
  \label{eq:dirty-grand-price}
\end{equation}

Fix $0<\lambda<t/2$ and let
\[
  L=\{(\{g\},\lambda):g=2,\ldots,m\}.
\]
Consider the star profile
\begin{equation}
  \theta^\star=\{H_h,(\{1\},a)\}\cup L.
  \label{eq:no-iig-star}
\end{equation}
By \eqref{eq:reverse-slack} and \eqref{eq:dirty-grand-price}, the grand bidder
faces a price strictly above $h$ and loses.

We next show that the high singleton must win.  Suppose instead that it loses
and that a low singleton $(\{g\},\lambda)$ wins.  Let
$L'=L\setminus\{(\{g\},\lambda)\}$.  Half-price persistence gives
\[
  c_{\{H_h\}\cup L'}(\{g\})\ge t.
\]
Against the external profile $\{H_h\}\cup L'$, let one physical bidder submit
the two reports $(\{g\},\lambda)$ and $(\{1\},a)$.  At $\theta^\star$, the
first wins and the second loses.  The false-name family covers $\{g\}$ and
pays at most $\lambda<t$,
contradicting Lemma~\ref{lem:det-cover}.  Thus no low singleton can win when
the high singleton loses.  Since the grand bidder already loses, this would
leave zero welfare on a positive-value profile, contradicting the assumed
positive ratio.  Hence $(\{1\},a)$ wins.

Finally, Lemma~\ref{lem:no-iig-half} implies that none of the low singletons
co-wins with $(\{1\},a)$.  Therefore
\[
  \SW_{\mathcal M}(\theta^\star)=a,
  \qquad
  \OPT(\theta^\star)\ge a+(m-1)\lambda,
\]
and
\[
  \rho(\mathcal M)
  \le\frac{a}{a+(m-1)\lambda}.
\]
Let $a\downarrow\bar t$ and $\lambda\uparrow t/2$.  Since
$\bar t\le(1+\eps)t$,
\[
  \rho(\mathcal M)
  \le
  \frac{2(1+\eps)}{2(1+\eps)+m-1}.
\]
Letting $\eps\downarrow0$ gives $2/(m+1)$.  For $m=1$, the statement reduces
to the trivial bound $\rho\le1$.
\end{proof}

The bound is tight: the Adaptive Reserve Price mechanism of
\cite{iwasaki2010worst} attains $2/(m+1)$.

\section{A randomized mechanism}
\label{sec:randomized}

Randomization removes the linear barrier of
Section~\ref{sec:deterministic}. The mechanism of this section is
a greedy auction with two sources of randomness. Every good is
deleted independently with a small probability, so that a report
is ranked by the value it is expected to realize, and winners
face a common random reserve, a uniform fraction of the largest
rejected weight. The deletion probability is governed by a
\emph{scale}.

We proceed in three steps. For every fixed scale, the mechanism
is universally truthful and false-name-proof in expectation
(Sections~\ref{sec:rand-greedy} and~\ref{sec:rand-scale}). Both
properties are therefore inherited by every report-independent
mixture of scales (Section~\ref{sec:rand-mixture}). Only the
welfare guarantee uses the particular mixture
(Section~\ref{sec:rand-welfare}).

\subsection{The mechanism}
\label{sec:rand-mechanism}

\paragraph{Greedy notation.}
A \emph{weighted report} \((S_i,w_i)\) consists of a target and a
weight \(w_i>0\). For an indexed collection \(T\) of weighted
reports, let \(G(T)\) be the packing obtained by processing
reports in nonincreasing order of weight and accepting a report
whenever its target is disjoint from the previously accepted
targets. Ties are resolved by a priority order \(\omega\) on the
identities, independent of the reported values and targets, and
every subprofile inherits the induced order. Define
\[
C(T)=\max_{i\in T\setminus G(T)}w_i,
\]
the largest rejected weight, where the maximum of an empty set is
zero. For a weighted profile \(E\) and a nonempty bundle
\(R\subseteq M\), define
\[
\kappa(R;E)=
\max\bigl\{
w_j:j\in G(E),\
S_j\cap R\ne\varnothing
\bigr\},
\qquad
B(R;E)=\max\{C(E),\kappa(R;E)\},
\]
again using zero for an empty maximum. The \emph{greedy critical
weight} of \(i\in G(T)\) is
\(\kappa_i(T)=\kappa(S_i;T\setminus\{i\})\).

\paragraph{Scales.}
Let
\[
\mathcal K=\bigl\{1,2,4,\ldots,2^{\lfloor\log_2m\rfloor}\bigr\},
\qquad
|\mathcal K|=1+\lfloor\log_2m\rfloor.
\]
A \emph{scale} is an element \(k\in\mathcal K\cup\{\infty\}\). At
scale \(k\), every good survives with probability
\[
p_k=1-\frac{1}{2k},
\qquad
p_\infty=1,
\]
so that a bundle \(R\) survives with probability
\(q_k(R)=p_k^{|R|}\).

\begin{definition}[Fixed-scale mechanism \(\mathcal A_k\)]
\label{def:Ak}
Fix a scale \(k\). Reports with value zero are discarded. Given
the remaining reports \((S_i,b_i)\), proceed as follows.
\begin{enumerate}
\item \emph{Weights.} Give report \(i\) the weight
\(w_i=q_k(S_i)b_i\), and let \(T\) be the resulting weighted
profile.
\item \emph{Greedy.} Draw a uniformly random priority order
\(\omega\) on the participating identities, independently of the
reports, and compute
\[
G=G(T),\qquad
C=C(T),\qquad
\kappa_i=\kappa_i(T)\quad(i\in G).
\]
\item \emph{Reserve and deletions.} Draw
\(\xi\sim\operatorname{Unif}[0,1]\). Independently, form
\(\widetilde M\subseteq M\) by retaining each good independently
with probability \(p_k\).
\item \emph{Allocation and payments.} Bidder \(i\) receives
\(S_i\) precisely when
\[
i\in G,\qquad w_i\ge\xi C,\qquad S_i\subseteq\widetilde M,
\]
and then pays
\begin{equation}\label{eq:actual-payment}
P_i=\frac{\max\{\kappa_i,\xi C\}}{q_k(S_i)}.
\end{equation}
Every other bidder receives nothing and pays zero.
\end{enumerate}
\end{definition}

In words, \(\mathcal A_k\) ranks every report by its expected
surviving value \(q_k(S_i)b_i\), runs greedy, and serves the
greedy winners whose weight clears the common reserve \(\xi C\)
and whose target survives. A served bidder pays the smallest bid
at which it would still have been served
(Lemma~\ref{lem:threshold}). At scale \(\infty\), no good is
deleted and weights equal bids.

\begin{definition}[The mechanism \(\mathcal M\)]
\label{def:M}
With probability \(1/2\), run \(\mathcal A_\infty\). With
probability \(1/2\), choose a scale \(k\) uniformly from
\(\mathcal K\) and run \(\mathcal A_k\).
\end{definition}

\begin{theorem}\label{thm:main}
For every \(m\ge1\), the mechanism \(\mathcal M\) of
Definition~\ref{def:M} for single-minded bidders is universally
truthful, false-name-proof in expectation, ex-post individually
rational, and makes no positive transfers. For every type
profile \(\theta\),
\[
\E[\SW_{\mathcal M}(\theta)]
\ge
\frac{\OPT(\theta)}{
4\sqrt{2m\ln(em)\bigl(1+\lfloor\log_2m\rfloor\bigr)}+8/e}.
\]
In particular, its worst-case welfare approximation ratio
\(1/\rho(\mathcal M)\) is
\[
O(\sqrt m\log(m+1))=\widetilde O(\sqrt m).
\]

The mechanism is anonymous in utilities and symmetric across
goods, and it uses polynomially many arithmetic operations. No
independence-of-irrelevant-goods condition is imposed.
\end{theorem}

Theorem~\ref{thm:main} combines
Corollary~\ref{cor:mixture}, Proposition~\ref{prop:welfare}, and
Proposition~\ref{prop:implementation}, which are proved in the
remainder of this section.

\subsection{Two facts about greedy packing}
\label{sec:rand-greedy}

\begin{lemma}[Greedy normalizer]
\label{lem:greedy}
Let \(T\) and \(E\) be weighted profiles.
\begin{enumerate}
\item[(a)] We have
\begin{equation}\label{eq:pair}
C(T)=
\max_{\substack{i,j\in T,\ i\ne j\\
                 S_i\cap S_j\ne\varnothing}}
\min\{w_i,w_j\}.
\end{equation}
In particular, \(C(T)\) does not depend on the priority order and
cannot decrease when reports are added.
\item[(b)] Add to \(E\) a report \((R,w)\) under a new identity.
It is accepted by greedy if \(w>\kappa(R;E)\) and rejected if
\(w<\kappa(R;E)\). Whenever it is accepted,
\begin{equation}\label{eq:normalizer}
C\bigl(E\cup\{(R,w)\}\bigr)=B(R;E).
\end{equation}
\end{enumerate}
\end{lemma}

\begin{proof}
(a) At least one member of every conflicting pair is rejected,
and the rejected member has weight at most \(C(T)\); hence the
minimum weight of the pair is at most \(C(T)\). Conversely, every
rejected report is blocked by an accepted report of at least its
weight, and the minimum weight of this pair is the rejected
weight. This gives the reverse inequality.

(b) Reports processed before \((R,w)\) are treated exactly as in
the greedy run on \(E\). The new report is therefore accepted
precisely when no report of \(G(E)\) processed before it has a
target intersecting \(R\). If \(w>\kappa(R;E)\), every such
report has smaller weight and is processed later. If
\(w<\kappa(R;E)\), a report of \(G(E)\) of weight
\(\kappa(R;E)\) whose target intersects \(R\) is processed
earlier.

Suppose now that the new report is accepted, so that
\(w\ge\kappa(R;E)\). A report of \(E\) that conflicts with it is
either accepted in \(G(E)\), and therefore has weight at most
\(\kappa(R;E)\), or rejected there, and therefore has weight at
most \(C(E)\). Conflicting pairs inside \(E\) have minimum weight
at most \(C(E)\). Equation~\eqref{eq:pair} gives the upper bound
\(C(E\cup\{(R,w)\})\le B(R;E)\). The lower bound follows from
monotonicity and, when \(\kappa(R;E)>0\), from the conflicting
report of \(G(E)\) of weight \(\kappa(R;E)\le w\).
\end{proof}

\subsection{Incentives at a fixed scale}
\label{sec:rand-scale}

Throughout this subsection, the scale
\(k\in\mathcal K\cup\{\infty\}\) is fixed, and a profile of
reports is identified with its weighted profile at scale \(k\).

\begin{lemma}[Realized critical prices]
\label{lem:threshold}
Fix a realization of \((\omega,\xi,\widetilde M)\) and an
external profile \(E\). For a nonempty bundle \(R\subseteq M\),
let
\[
c_E(R)=
\begin{cases}
\dfrac{\max\{\kappa(R;E),\,\xi B(R;E)\}}{q_k(R)},
&R\subseteq\widetilde M,\\[3mm]
+\infty,&\text{otherwise.}
\end{cases}
\]
A bidder reporting \((R,b)\) against \(E\) is served and pays
\(c_E(R)\) if \(b>c_E(R)\), and is not served if \(b<c_E(R)\).
If \(b=c_E(R)\), it is either served at payment \(b\) or not
served. Moreover, \(c_E(R)\le c_E(R')\) whenever
\(R\subseteq R'\).
\end{lemma}

\begin{proof}
A report whose target does not survive is never served. Assume
\(R\subseteq\widetilde M\) and let \(w=q_k(R)b\). By
Lemma~\ref{lem:greedy}(b), the report is accepted if
\(w>\kappa(R;E)\) and rejected if \(w<\kappa(R;E)\); whenever it
is accepted, its greedy critical weight is \(\kappa(R;E)\) and
the normalizer of the full profile is \(B(R;E)\). An accepted
report is thus served precisely when \(w\ge\xi B(R;E)\), and it
then pays
\(\max\{\kappa(R;E),\xi B(R;E)\}/q_k(R)=c_E(R)\)
by~\eqref{eq:actual-payment}. Consequently, the report is served
if \(w>\max\{\kappa(R;E),\xi B(R;E)\}\), it is not served if
\(w\) is smaller than this maximum, and at equality it pays
\(c_E(R)=b\) if it is served.

For monotonicity, let \(R\subseteq R'\). If \(R'\) does not
survive, there is nothing to prove. Otherwise \(R\) survives as
well, and
\[
\kappa(R;E)\le\kappa(R';E),
\qquad
B(R;E)\le B(R';E),
\qquad
q_k(R)\ge q_k(R')>0.
\qedhere
\]
\end{proof}

\begin{proposition}[Ex-post properties]
\label{prop:universal}
The mechanism \(\mathcal A_k\) is ex-post feasible, universally
truthful, ex-post individually rational, and makes no positive
transfers.
\end{proposition}

\begin{proof}
The served bidders form a subset of a greedy packing, and
payments are nonnegative. Fix a realization of the randomness,
an external profile \(E\), and a bidder of type \((S,v)\). By
Lemma~\ref{lem:threshold}, its truthful realized utility is
\(\max\{0,v-c_E(S)\}\ge0\). Consider any other report \((R,b)\).
If it is not served, or if \(S\not\subseteq R\), the realized
utility is at most zero. If it is served and \(S\subseteq R\),
the realized utility is
\[
v-c_E(R)\le v-c_E(S).
\]
In every case, the deviation does not improve upon the truthful
realized utility.
\end{proof}

We next compute the quantities that enter the false-name
argument. For \(z\ge0\) and \(0\le\kappa\le B\), define
\begin{equation}\label{eq:utility}
\Psi(z;\kappa,B)
=\E_\xi\bigl[\max\{0,\,z-\max\{\kappa,\xi B\}\}\bigr].
\end{equation}

\begin{lemma}[Expected payments and truthful utility]
\label{lem:interim}
Fix the priority order \(\omega\), and take expectations over
\(\xi\) and \(\widetilde M\).
\begin{enumerate}
\item[(a)] Let \(i\in G\). Then
\(\kappa_i\le\min\{w_i,C\}\). Bidder \(i\) is served with
probability \(q_k(S_i)r_i\) and has expected payment \(\bar P_i\),
where, if \(C>0\),
\begin{equation}\label{eq:retention-payment}
r_i=\min\left\{1,\frac{w_i}{C}\right\},
\qquad
\bar P_i=\frac{\min\{w_i,C\}^2+\kappa_i^2}{2C},
\end{equation}
and \(r_i=1\), \(\bar P_i=0\) if \(C=0\).
\item[(b)] Consider a truthful bidder of type \((S,v)\) facing an
external profile \(E\), and let
\[
z=q_k(S)v,
\qquad
\kappa=\kappa(S;E),
\qquad
B=B(S;E).
\]
Its expected utility equals \(\Psi(z;\kappa,B)\). Moreover,
\(\Psi(z;\kappa,B)=0\) if \(z<\kappa\), and
\[
\Psi(z;\kappa,B)
=\int_\kappa^z\min\left\{\frac{\zeta}{B},1\right\}\,d\zeta
\qquad(z\ge\kappa),
\]
with integrand one if \(B=0\). In particular, for fixed \(z\) and
\(\kappa\), this utility is nonincreasing in \(B\ge\kappa\), and
for \(B>0\),
\[
\Psi(z;\kappa,B)=
\begin{cases}
0,&z<\kappa,\\[1mm]
\dfrac{z^2-\kappa^2}{2B},&\kappa\le z\le B,\\[3mm]
z-\dfrac{B^2+\kappa^2}{2B},&z\ge B.
\end{cases}
\]
\end{enumerate}
\end{lemma}

\begin{proof}
(a) Acceptance gives \(\kappa_i\le w_i\), and
Lemma~\ref{lem:greedy}(b), applied to \(T\setminus\{i\}\), gives
\(C=B(S_i;T\setminus\{i\})\ge\kappa_i\). The target \(S_i\)
survives with probability \(q_k(S_i)\), independently of \(\xi\),
and \(w_i\ge\xi C\) holds with probability \(r_i\).
By~\eqref{eq:actual-payment}, the factor \(q_k(S_i)\) cancels in
the expected payment, which equals
\(\E_\xi[\mathbf{1}\{\xi C\le w_i\}\max\{\kappa_i,\xi C\}]\).
This is zero if \(C=0\). If \(C>0\), it equals
\[
\int_0^{\kappa_i/C}\kappa_i\,d\xi
+\int_{\kappa_i/C}^{r_i}\xi C\,d\xi
=\frac{\kappa_i^2}{C}
+\frac{\min\{w_i,C\}^2-\kappa_i^2}{2C}
=\bar P_i.
\]

(b) By Lemma~\ref{lem:threshold}, the truthful realized utility
is \(\max\{0,v-c_E(S)\}\), which vanishes unless \(S\) survives.
Since \(S\) survives with probability \(q_k(S)\), independently
of \(\xi\), the expected utility is
\[
q_k(S)\,
\E_\xi\!\left[
\max\left\{0,\,
v-\frac{\max\{\kappa,\xi B\}}{q_k(S)}\right\}
\right]
=\Psi(z;\kappa,B).
\]
If \(z<\kappa\), the integrand in~\eqref{eq:utility} is zero. If
\(z\ge\kappa\), then
\[
\max\{0,\,z-\max\{\kappa,\xi B\}\}
=\int_\kappa^z\mathbf{1}\{\xi B\le\zeta\}\,d\zeta,
\]
and taking the expectation over \(\xi\) gives the integral
formula. The remaining claims follow by direct computation.
\end{proof}

The next lemma is the combinatorial core of the false-name
argument. It concerns greedy packing only.

\begin{lemma}[Payment certificate]
\label{lem:certificate}
Let \(E\) be a weighted profile, let \(S\subseteq M\) be a
nonempty bundle, and put \(\kappa=\kappa(S;E)\) and
\(B=B(S;E)\). Let \(T'\) be obtained from \(E\) by adding
finitely many weighted reports under new identities \(A\), with
a priority order extending that of \(E\). Let \(F\subseteq A\)
be the set of new reports accepted by \(G(T')\), write
\(C'=C(T')\), and suppose that the targets of \(F\) cover \(S\).
Define
\[
F_S=\{j\in F:S_j\cap S\ne\varnothing\},
\qquad
\alpha=\min_{j\in F_S}\min\{w_j,C'\}.
\]
Then
\begin{equation}\label{eq:C-comparison}
C'\ge B,
\end{equation}
and
\begin{equation}\label{eq:certificate}
\sum_{j\in F}
\left(
\min\{w_j,C'\}^2+\kappa_j(T')^2
\right)
\ge\alpha^2+\kappa^2.
\end{equation}
\end{lemma}

\begin{proof}
By Lemma~\ref{lem:greedy}(a), \(C'\ge C(E)\). If \(\kappa>0\),
fix a report of \(G(E)\) of weight \(\kappa\) whose target
intersects \(S\), and call it the \emph{blocker}. The blocker
cannot be accepted in \(G(T')\), because the targets of \(F\)
cover \(S\) and a greedy packing is disjoint. It is therefore
rejected in \(T'\), implying \(C'\ge\kappa\). This
proves~\eqref{eq:C-comparison}.

Choose \(j_0\in F_S\) attaining the minimum defining \(\alpha\).
Its capped-weight term supplies \(\alpha^2\). If \(\kappa=0\),
the claim follows immediately. If some other report of \(F\) has
weight at least \(\kappa\), its capped-weight term supplies
\(\kappa^2\), because \(C'\ge\kappa\).

It remains to consider the case in which \(\kappa>0\) and every
report of \(F\) other than \(j_0\) has weight below \(\kappa\).
We show that \(\kappa_{j_0}(T')\ge\kappa\), so that the
critical-weight term of \(j_0\) supplies \(\kappa^2\). Suppose,
to the contrary, that \(\kappa_{j_0}(T')<\kappa\): in the profile
\(T'\setminus\{j_0\}\), no accepted report of weight at least
\(\kappa\) has a target intersecting \(S_{j_0}\).

Call the reports of weight at least \(\kappa\) the \emph{prefix};
greedy processes them first. We claim that \(G(T')\) and
\(G(T'\setminus\{j_0\})\) accept the same prefix reports other
than \(j_0\). This is automatic if \(w_{j_0}<\kappa\). Otherwise,
scan the prefix in greedy order. The two runs coincide before
\(j_0\), and \(j_0\) is accepted in \(G(T')\). A later prefix
report accepted in \(G(T'\setminus\{j_0\})\) has a target
disjoint from \(S_{j_0}\) by assumption, so the presence of
\(j_0\) never changes a decision.

By the claim, a new report in the prefix that is accepted in
\(G(T'\setminus\{j_0\})\) belongs to \(F\setminus\{j_0\}\) and
has weight at least \(\kappa\); no such report exists. Hence all
new reports in the prefix are rejected in
\(G(T'\setminus\{j_0\})\). Deleting rejected reports never
changes a greedy outcome, so the prefix reports accepted in
\(G(T'\setminus\{j_0\})\) are exactly the prefix reports of
\(G(E)\). In particular, the blocker is accepted in
\(G(T'\setminus\{j_0\})\) and, by the claim, in \(G(T')\). This
contradicts the first paragraph of the proof.
\end{proof}

\begin{proposition}[False-name-proofness at a fixed scale]
\label{prop:fnp-scale}
The mechanism \(\mathcal A_k\) is false-name-proof in
expectation.
\end{proposition}

\begin{proof}
Fix a bidder of type \((S,v)\), the external profile \(E\), and
the priority order among the external identities. Let \(z\),
\(\kappa\), and \(B\) be as in Lemma~\ref{lem:interim}(b). The
truthful expected utility is \(\Psi(z;\kappa,B)\), whatever the
bidder's own priority.

Replace the truthful report by any finite family
\(\hat\theta_A\) of alias reports, with arbitrary targets,
values, and priorities. Let \(T'\) be the resulting weighted
profile, let \(F\subseteq A\) be the set of aliases accepted by
\(G(T')\), and write \(C'=C(T')\). Abstention gives zero utility.
If the targets of \(F\) do not cover \(S\), the deviation
produces no value and has nonnegative payments. In both cases it
cannot improve upon truthful utility. Suppose henceforth that
the targets of \(F\) cover \(S\), and let \(F_S\) and \(\alpha\)
be as in Lemma~\ref{lem:certificate}.

Assume first that \(C'>0\). Every alias in \(F_S\) is necessary
to supply \(S\), because the accepted targets are disjoint.
These aliases clear the common reserve \(\xi C'\) simultaneously
with probability \(\alpha/C'\). Their targets must also all
survive, which occurs with probability \(q_k(Q)\le q_k(S)\),
where
\[
Q=\bigcup_{j\in F_S}S_j\supseteq S.
\]
Hence the deviation's expected value is at most
\begin{equation}\label{eq:alias-value}
v\,q_k(S)\frac{\alpha}{C'}=\frac{z\alpha}{C'}.
\end{equation}
The bidder pays for all accepted aliases, including helpers
whose targets are irrelevant to \(S\). By
Lemma~\ref{lem:interim}(a) and~\eqref{eq:certificate}, its total
expected payment is
\[
\sum_{j\in F}\bar P_j
=\sum_{j\in F}
\frac{\min\{w_j,C'\}^2+\kappa_j(T')^2}{2C'}
\ge\frac{\alpha^2+\kappa^2}{2C'}.
\]
Together with~\eqref{eq:alias-value}, this gives
\begin{equation}\label{eq:attack}
\E[u_{\mathrm{aliases}}]
\le
\frac{z\alpha-(\alpha^2+\kappa^2)/2}{C'},
\qquad
0\le\alpha\le C'.
\end{equation}
If \(z<\kappa\), the right-hand side is at most
\((z-\kappa)\alpha/C'\le0\), whereas truthful utility is zero.
If \(z\ge\kappa\), its maximum over \(0\le\alpha\le C'\) is
exactly \(\Psi(z;\kappa,C')\).
Using~\eqref{eq:C-comparison} and monotonicity in the last
argument,
\[
\E[u_{\mathrm{aliases}}]
\le
\Psi(z;\kappa,C')
\le
\Psi(z;\kappa,B)
=
\E[u_{\mathrm{truthful}}].
\]

If \(C'=0\), then \(B=0\) by~\eqref{eq:C-comparison}, all
payments are zero, and the deviation succeeds with probability
at most \(q_k(S)\). Truthful utility is
\(\Psi(z;0,0)=z=q_k(S)v\), so the desired inequality holds in
this case as well.

The argument is conditional on the external priorities and
permits arbitrary priorities for the aliases. It therefore
remains valid after averaging over the uniformly random priority
order.
\end{proof}

\subsection{Mixtures of scales}
\label{sec:rand-mixture}

\begin{corollary}\label{cor:mixture}
Every mechanism that draws a scale
\(k\in\mathcal K\cup\{\infty\}\) from a fixed distribution,
independently of the reports, and then runs \(\mathcal A_k\) is
ex-post feasible, universally truthful, ex-post individually
rational, false-name-proof in expectation, and makes no positive
transfers. In particular, this holds for \(\mathcal M\).
\end{corollary}

\begin{proof}
By Proposition~\ref{prop:universal}, the ex-post properties hold
for every realization of the scale. Because the scale is drawn
independently of the reports, the expected utility of any family
of reports is the average over the scale of its expected utility
in \(\mathcal A_k\). The inequality of
Proposition~\ref{prop:fnp-scale} is preserved by this average.
\end{proof}

\subsection{Welfare of the mixture}
\label{sec:rand-welfare}

\begin{proposition}[Welfare guarantee]
\label{prop:welfare}
For every type profile \(\theta\),
\[
\E[\SW_{\mathcal M}(\theta)]
\ge
\frac{\OPT(\theta)}{
4\sqrt{2m\ln(em)\,|\mathcal K|}+8/e}.
\]
\end{proposition}

\begin{proof}
Assume a truthful profile \(\theta\) and put \(V=\max_i v_i\).
If \(V=0\), the claim is immediate.

For a scale \(k\in\mathcal K\cup\{\infty\}\), let \(\Phi_k\)
denote the expected welfare of \(\mathcal A_k\), including
priority randomization. For a fixed priority order \(\omega\),
let \(G_k\) and \(C_k\) denote the greedy packing and the
normalizer \(C(T)\) at scale \(k\), and let
\(\varphi_{k,\omega}\) denote the expected welfare of
\(\mathcal A_k\) over \(\xi\) and \(\widetilde M\). By
Lemma~\ref{lem:interim}(a),
\[
\varphi_{k,\omega}
=
\sum_{i\in G_k}
w_i\min\left\{1,\frac{w_i}{C_k}\right\},
\]
with the minimum replaced by one if \(C_k=0\).
Since \(w_i\le V\) and \(C_k\le V\),
\begin{equation}\label{eq:squares}
\sum_{i\in G_k}w_i^2\le V\varphi_{k,\omega}.
\end{equation}

\paragraph{The scale \(\infty\).}
At scale \(\infty\), weights equal values and no good is
deleted. The report processed first by greedy has value \(V\),
is accepted, and has weight at least \(C_\infty\); it is
therefore served for every \(\xi\). Hence
\begin{equation}\label{eq:top}
\Phi_\infty\ge V.
\end{equation}

\paragraph{The scales in \(\mathcal K\).}
Fix an optimal packing \(N^\star\subseteq N\), that is, a set of
bidders with pairwise disjoint targets, positive values, and
total value \(\OPT(\theta)\). For each \(k\in\mathcal K\), let
\[
N_k=\{i\in N^\star:k\le |S_i|<2k\},
\]
and define
\[
W_k=\sum_{i\in N_k}v_i,\qquad
n_k=|N_k|,\qquad
m_k=\sum_{i\in N_k}|S_i|.
\]
Every target has between \(1\) and \(m\) goods, so the classes
\(N_k\) partition \(N^\star\). Disjointness gives
\begin{equation}\label{eq:mass}
k\,n_k\le m_k,\qquad
\sum_{k\in\mathcal K}m_k\le m,\qquad
\sum_{k\in\mathcal K}n_k\le m.
\end{equation}
For \(i\in N_k\), we have \(|S_i|<2k\), and
\((1-1/(2k))^{2k}\) is nondecreasing in \(k\ge1\), so
\[
q_k(S_i)
\ge
\left(1-\frac{1}{2k}\right)^{2k}
\ge\left(1-\frac12\right)^{2}
=\frac14,
\]
and hence \(v_i\le4w_i\).

For the fixed priority order, assign each bidder of \(N_k\) to
the first greedy winner that blocks it, or to itself if it is
accepted. Let \(\nu_j\) be the number assigned to greedy winner
\(j\in G_k\). Greedy ordering and disjointness of \(N^\star\)
imply
\[
W_k\le4\sum_{j\in G_k}\nu_jw_j,\qquad
\nu_j\le|S_j|,\qquad
\sum_{j\in G_k}\nu_j=n_k.
\]

Set \(d=\ln(em)\). For winners with \(|S_j|\le2kd\),
\[
\sum_{\substack{j\in G_k\\|S_j|\le2kd}}\nu_j^2
\le 2kd\,n_k
\le 2d\,m_k.
\]
By Cauchy--Schwarz and~\eqref{eq:squares},
\[
\sum_{\substack{j\in G_k\\|S_j|\le2kd}}\nu_jw_j
\le
\sqrt{2d\,m_kV\varphi_{k,\omega}}.
\]

For winners with \(|S_j|>2kd\),
\[
q_k(S_j)
\le
e^{-|S_j|/(2k)}
<
\frac1{em},
\]
so
\[
\sum_{\substack{j\in G_k\\|S_j|>2kd}}\nu_jw_j
\le
\frac{Vn_k}{em}.
\]
Combining these inequalities and averaging over \(\omega\),
using concavity of the square root and
\(\Phi_k=\E_\omega[\varphi_{k,\omega}]\), gives
\begin{equation}\label{eq:class}
W_k
\le
4\sqrt{2d\,m_kV\Phi_k}
+
\frac{4Vn_k}{em}.
\end{equation}

Sum~\eqref{eq:class} over \(k\in\mathcal K\) and apply
Cauchy--Schwarz and~\eqref{eq:mass}:
\begin{equation}\label{eq:opt}
\begin{aligned}
\OPT(\theta)
&\le
4\sqrt{2dV}\sum_{k\in\mathcal K}\sqrt{m_k\Phi_k}
+\frac{4V}{e}\\
&\le
4\sqrt{2mdV\sum_{k\in\mathcal K}\Phi_k}
+\frac{4V}{e}.
\end{aligned}
\end{equation}

\paragraph{The mixture.}
By Definition~\ref{def:M} and~\eqref{eq:top}, the expected
welfare \(\Phi\) of \(\mathcal M\) satisfies
\[
\Phi
=\frac12\Phi_\infty
+\frac1{2|\mathcal K|}\sum_{k\in\mathcal K}\Phi_k
\ge\frac12V
+\frac1{2|\mathcal K|}\sum_{k\in\mathcal K}\Phi_k.
\]
The arithmetic--geometric mean inequality yields
\[
V\sum_{k\in\mathcal K}\Phi_k\le|\mathcal K|\,\Phi^2,
\qquad
V\le2\Phi.
\]
Substituting into~\eqref{eq:opt} proves
\[
\OPT(\theta)
\le
\left(
4\sqrt{2m\ln(em)\,|\mathcal K|}+\frac8e
\right)\Phi,
\]
which is the claimed approximation guarantee.
\end{proof}

\subsection{Implementation and symmetry}
\label{sec:rand-implementation}

\begin{proposition}\label{prop:implementation}
The mechanism \(\mathcal M\) is anonymous in utilities and
symmetric across goods. For \(n\) submitted reports, it uses
\(O(n^2\log n+n^2m)\) arithmetic operations.
\end{proposition}

\begin{proof}
The priority order is uniformly random on the participating
identities and independent of their reports, so relabelling
bidders does not change the distribution of outcomes. All other
operations depend on goods only through cardinalities,
intersections, and independent identical survival
probabilities. The mechanism is therefore anonymous in utilities
and symmetric across goods.

Only one scale is executed. One greedy run takes
\(O(n\log n+nm)\) elementary operations with explicit set
representations, and computing all needed critical weights by
deletion runs takes \(O(n^2\log n+n^2m)\) operations.
\end{proof}

\paragraph{Exact arithmetic.}
The payment~\eqref{eq:actual-payment} is a continuous random
variable whenever \(\xi C>\kappa_i\). If exact finite-precision
output is required, a served bidder can instead be charged its
conditional expected payment
\(\bar P_i/(q_k(S_i)r_i)\), with \(r_i\) and \(\bar P_i\) as
in~\eqref{eq:retention-payment}. The allocation and all expected
payments are unchanged, so the welfare guarantee,
truthfulness in expectation, and false-name-proofness in
expectation are preserved; since
\(\bar P_i\le w_ir_i\), this payment is at most \(b_i\), and the
mechanism remains ex-post individually rational. It is no longer
universally truthful, because a winner's payment then depends on
its own bid. For rational reports, \(q_k(S)\) has denominator
\((2k)^{|S|}\), so weights, service probabilities, and these
payments have polynomial bit length. The reserve lottery can be
implemented by sampling among the at most \(|G|+1\) intervals
determined by the rational thresholds \(r_i\), and standard
exact rational sampling has expected polynomial bit complexity.

\paragraph{Scope of the result.}
The false-name guarantee is in expectation, with the bidder
liable for the payments of all its aliases. The deletion step
does not impose an extension of independence of irrelevant
goods: adding a privately desired good changes the probability
that the entire target survives.

\bibliographystyle{ACM-Reference-Format}
\bibliography{sample-bibliography}
\end{document}
